\section{Discussion}
\subsection{Common-grid agreement and geographic heterogeneity}
Harmonization to the FCOVER target grid removes a direct native-pixel-versus-300 m comparison, but not differences in spectral response, timing, atmosphere, geometry, QA, or effective valid support \citep{claverie2018,zhang2018ndvi,shang2019,trevisiol2024}. The matched aggregation-order sensitivity preserved all target identities and the geographic interpretation, although it modified some sensor--AOI estimates and one preferred history; it does not establish either aggregation order as universally correct. MODIS had the lowest mean RMSE in this protocol, but its native support and aggregation choice may contribute; this is not a universal sensor-capability ranking \citep{woodcock1987,atkinson2000}. Geographic variation exceeded the sensor-mean separation. AOI-01 particularly shows why low absolute error must be read with target range, $R^2$, and coefficient identifiability. The four AOIs reveal conditionality but are designed contrasts, not a regional probability sample \citep{roberts2017,meyer2018,ploton2020}.

\subsection{Historical-window transfer}
No history was geographically invariant, and Rolling-Origin results did not show monotonic benefit from longer windows. The primary $\pm15$ d windows overlap; the non-overlapping assignment changed several selected histories and directions, but preserved the absence of a uniform longer-history benefit. Coefficient drift and the diagnostic variation are consistent with changing NDVI--FCOVER relations across place and history. Because a historical window changes years, blocks, pair counts, and their implicit weights together, these results do not identify the effect of adding a year alone. The practical implication is to assess the intended domain with block-grouped diagnostics and chronological forward evaluation rather than treating record length as a universal tuning direction \citep{tashman2000,bergmeir2012}.

\subsection{DPM and practical implications}
The DPM results show that empirical-endpoint sensitivity is geographically conditional rather than invariant. AOI-02 was comparatively stable, while AOI-03 had the largest absolute endpoint-driven RMSE differences; AOI-01 was not uniquely unstable, although its P1/P99 empirical high-NDVI endpoint proxy still produced large positive bias. Narrowing the quantile range did not rescue AOI-01, indicating that its extreme DPM error reflects restricted FCOVER support and the empirical-reference-range mismatch in addition to endpoint parameterization. The AOI-01 OLS fits retained modest skill against zero and historical-mean baselines, but their low absolute errors should continue to be read with the restricted target range and the documented clipping frequency. DPM remains useful when paired calibration is unavailable, but empirical endpoints are not measured endmembers and estimates remain sensitive to background, shadow, and two-component assumptions \citep{huete1988,montandon2008,jiapaer2011,yan2022dpm,mu2024}. Configuration choice should therefore be evaluated in its intended geographic and temporal setting.

Landsat Collection 2 Level-2 surface reflectance was already atmospherically corrected in the primary analysis. Additional \texttt{SR\_QA\_AEROSOL} screening was associated with configuration-specific error changes, including large changes after strict filtering in AOI-02, rather than evidence that aerosol alone caused those differences. Its support loss under strict rules reinforces the need to report the retained sample and residual spectral, atmospheric, and support differences alongside any aerosol-screened result.

\section{Limitations}
Copernicus FCOVER is an operational product reference, not field ground truth; field, UAV, or validated high-resolution reference data would be required for absolute FVC accuracy \citep{baret2013,camacho2013,fuster2020,copernicus2025validation,mu2015sampling,chen2016uav,olofsson2014}. The four designed AOIs are not plateau probability samples. A common 300 m grid does not equalize effective valid support: no exact valid-area weighting, overlap denominator, or minimum native-pixel threshold was retained. Matched-support aggregation-order results bound one processing effect but do not resolve residual cross-sensor support differences or limitations of mean aggregation. Primary nominal-date windows overlap; the unique-assignment sensitivity removed duplicate source use but also changed support. Native QA cannot make acquisition, atmospheric, spectral, or support conditions equivalent across sensors, and valid target subsets can consequently differ. FCOVER validity followed the implemented QFLAG/NOBS rules, but QFLAG classes and FCOVER reference-quality subsets were not evaluated. Landsat aerosol screening was evaluated, but strict modes can materially reduce support and do not remove residual atmospheric or spectral differences.

Historical windows differ in years, blocks, pair counts, and implicit weights, so causal effects of historical depth cannot be isolated. Paired target identities were identical within a block contrast, but the corresponding training samples were not. The reserve was recombined before full fitting; GroupKFold was not buffered extrapolation; the AOI-01 block scale is approximate; and residual spatial dependence may remain. No block-size, autocorrelation, permutation, or bootstrap sensitivity was performed. DPM endpoints remain empirical proxies, and selected DPM/OLS minima were retrospectively selected from different candidate and information conditions. AOI-01's near-zero distribution still constrains interpretation, although clipping frequency and naïve-baseline performance are now quantified.

\section{Conclusion}
Common-target-grid harmonization enabled sensor-specific NDVI--FCOVER agreement assessment without direct native-pixel-versus-target-grid comparison. Across four AOIs, agreement, coefficients, and preferred histories were geographically conditional, and longer historical-window configurations were not uniformly associated with lower forward error. These interpretations were robust to the tested aggregation-order, non-overlapping-composition, and Landsat aerosol-QA sensitivities, despite configuration-specific numerical changes. The practical implication is to evaluate the intended-domain history with both block-grouped stability diagnostics and chronological forward testing, rather than select a configuration by record length alone. Reporting local contrast magnitude and direction alongside inference makes this decision transparent and guards against over-reading pooled summaries. Its role is to guide configuration evaluation, not to rank sensors. The common target grid reduces one source of mismatch but does not equalize valid support after QA. These findings concern agreement with Copernicus FCOVER and do not establish absolute field-level FVC accuracy.
